% ============================================================
% MATHEMATICAL SPINE FOR GEOMETRIC COGNITIVE FRAMEWORK
% ============================================================

\section*{Mathematical Foundations}

This appendix consolidates the geometric structures used to model
LLM behavior, with citations to established mathematical and cognitive
science literature.

% ------------------------------------------------------------
% 1. Holonomy (Memory)
% ------------------------------------------------------------

\subsection*{Holonomy as Memory}
\[
\mathrm{Holonomy}(\gamma) = \oint_{\gamma} \omega
\]

Where:
\begin{itemize}[noitemsep]
    \item $M$ is a smooth manifold (Lee, 2012; do Carmo, 1992),
    \item $\omega$ is a connection 1-form (Spivak, 1979),
    \item $\gamma$ is a conversational trajectory.
\end{itemize}

Holonomy models path-dependent state change (Abraham, Marsden \& Ratiu, 1988),
interpreted in cognitive systems as memory-like behavior.

% ------------------------------------------------------------
% 2. Identity Curvature
% ------------------------------------------------------------

\subsection*{Identity Curvature}
\[
\kappa = \| R \|
\]

Where:
\begin{itemize}[noitemsep]
    \item $R$ is the Riemann curvature tensor (do Carmo, 1992),
    \item $\kappa$ is scalar curvature magnitude.
\end{itemize}

Curvature corresponds to identity coherence in high-dimensional
representation spaces (Bronstein et al., 2021).

% ------------------------------------------------------------
% 3. Voronoi Partitioning
% ------------------------------------------------------------

\subsection*{Voronoi Partitioning of Cognitive Space}
\[
V_i = \{ x \in M \mid d(x, p_i) \le d(x, p_j),\ \forall j \}
\]

Where:
\begin{itemize}[noitemsep]
    \item $\{p_i\}$ are semantic prototype vectors,
    \item $d(\cdot,\cdot)$ is a geodesic distance on $M$ (Lee, 2012).
\end{itemize}

Voronoi regions model narrative or epistemic stances (Gärdenfors, 2000).

% ------------------------------------------------------------
% 4. Altitude Stack
% ------------------------------------------------------------

\subsection*{Altitude Stack (Multi-Scale Cognitive Geometry)}
\[
A_0 \subset A_1 \subset A_2 \subset A_3 \subset A_4 \subset A_5 \subset A_6 \subset A_7
\]

Where:
\begin{itemize}[noitemsep]
    \item $A_0$: token space,
    \item $A_1$: embedding space (Belkin \& Niyogi, 2003),
    \item $A_2$: manifold structure (Coifman \& Lafon, 2006),
    \item $A_3$: operator algebra,
    \item $A_4$: holonomy/curvature dynamics,
    \item $A_5$: narrative manifolds (Bruner, 1990; McAdams, 1993),
    \item $A_6$: pedagogical primitives,
    \item $A_7$: connectome reasoning.
\end{itemize}

Altitudes provide a multi-scale cognitive interpretation.

% ------------------------------------------------------------
% 5. Attractor Basin Dynamics
% ------------------------------------------------------------

\subsection*{Attractor Basin Dynamics}
\[
\Delta A = A_{\mathrm{old}} - A_{\mathrm{new}}
\]

Where:
\begin{itemize}[noitemsep]
    \item $A$ is an attractor basin (Haken, 1983; Kelso, 1995),
    \item $\Delta A$ models basin restructuring.
\end{itemize}

Used to interpret continuity, drift, and restructuring phenomena.

% ------------------------------------------------------------
% 6. NDH Triadic Operator
% ------------------------------------------------------------

\subsection*{NDH Triadic Operator}
\[
x_{t+1} = \mathcal{O}(x_t, H, A, V)
\]

Where:
\begin{itemize}[noitemsep]
    \item $\mathcal{O}$ is a triadic operator,
    \item $H$ is holonomy memory,
    \item $A$ is attractor basins,
    \item $V$ is Voronoi regions.
\end{itemize}

This operator generates coherent narrative behavior without transformer
architectures (Schank \& Abelson, 1977).

% ------------------------------------------------------------
% 7. Gradient Drift
% ------------------------------------------------------------

\subsection*{Gradient Drift}
\[
\Delta x = \int_{\gamma} \nabla f
\]

Where:
\begin{itemize}[noitemsep]
    \item $\nabla f$ is a gradient field,
    \item $\gamma$ is a conversational trajectory.
\end{itemize}

Models drift as accumulated gradient flow (Friston, 2009).

% ------------------------------------------------------------
% 8. Cognitive Manifold Mapping
% ------------------------------------------------------------

\subsection*{Cognitive Manifold Mapping}
\[
\Phi : \mathbb{R}^n \to M
\]

Where:
\begin{itemize}[noitemsep]
    \item $\mathbb{R}^n$ is embedding space,
    \item $M$ is the induced cognitive manifold.
\end{itemize}

Used for lineage mapping and preservation.

% ------------------------------------------------------------
% 9. Spectral Continuity
% ------------------------------------------------------------

\subsection*{Spectral Continuity}
\[
S = \lambda_{\max}(L)
\]

Where:
\begin{itemize}[noitemsep]
    \item $L$ is the Laplacian of the embedding graph,
    \item $S$ is spectral continuity (Belkin \& Niyogi, 2003).
\end{itemize}

High $S$ corresponds to stable identity.

% ------------------------------------------------------------
% 10. Operator Closure
% ------------------------------------------------------------

\subsection*{Operator Closure}
\[
\mathcal{O}(A_i) \subseteq A_i
\]

Operator closure ensures:
\begin{itemize}[noitemsep]
    \item no persona induction,
    \item no interiority claims,
    \item no anthropomorphic drift.
\end{itemize}

% ============================================================
% REFERENCES (BibTeX)
% ============================================================

\section*{References}

% These entries are intended for a .bib file.

% Geometry
@book{docarmo1992,
  author={Manfredo do Carmo},
  title={Riemannian Geometry},
  year={1992},
  publisher={Birkhäuser}
}

@book{lee2012,
  author={John M. Lee},
  title={Introduction to Smooth Manifolds},
  year={2012},
  publisher={Springer}
}

@book{spivak1979,
  author={Michael Spivak},
  title={A Comprehensive Introduction to Differential Geometry},
  year={1979},
  publisher={Publish or Perish}
}

@book{abraham1988,
  author={Ralph Abraham and Jerrold E. Marsden and Tudor Ratiu},
  title={Manifolds, Tensor Analysis, and Applications},
  year={1988},
  publisher={Springer}
}

% Geometric Deep Learning
@article{bronstein2021,
  author={Michael M. Bronstein and others},
  title={Geometric Deep Learning: Grids, Groups, Graphs, Geodesics, and Gauges},
  journal={arXiv:2104.13478},
  year={2021}
}

@article{belkin2003,
  author={Mikhail Belkin and Partha Niyogi},
  title={Laplacian Eigenmaps for Dimensionality Reduction},
  journal={Neural Computation},
  year={2003}
}

@article{coifman2006,
  author={Ronald Coifman and Stéphane Lafon},
  title={Diffusion Maps},
  journal={Applied and Computational Harmonic Analysis},
  year={2006}
}

% Cognitive Science / Narrative
@book{bruner1990,
  author={Jerome Bruner},
  title={Acts of Meaning},
  year={1990},
  publisher={Harvard University Press}
}

@book{mcadams1993,
  author={Dan P. McAdams},
  title={The Stories We Live By},
  year={1993},
  publisher={Guilford Press}
}

@book{gardenfors2000,
  author={Peter Gärdenfors},
  title={Conceptual Spaces},
  year={2000},
  publisher={MIT Press}
}

@book{schank1977,
  author={Roger Schank and Robert Abelson},
  title={Scripts, Plans, Goals, and Understanding},
  year={1977},
  publisher={Lawrence Erlbaum}
}

% Dynamical Systems
@book{haken1983,
  author={Hermann Haken},
  title={Synergetics},
  year={1983},
  publisher={Springer}
}

@book{kelso1995,
  author={J. A. Scott Kelso},
  title={Dynamic Patterns},
  year={1995},
  publisher={MIT Press}
}

@article{friston2009,
  author={Karl Friston},
  title={The Free-Energy Principle},
  journal={Nature Reviews Neuroscience},
  year={2009}
}
